% ── Automatizace, Průmysl 4.0 a generativní AI ──────────────────────────────

Diskuse o~vlivu technologického pokroku na~trh práce má~v~kontextu sociálního
dialogu dvojí význam: kvantifikuje očekávaný objem strukturální fluktuace
zaměstnanosti a~zároveň určuje, ve~kterých segmentech bude poptávka po~koordinovaném
přizpůsobení (rekvalifikace, mzdové kompenzace, tarifní pravidla pro~nové formy
práce) nejvyšší. Aktuální odborná literatura nabízí překvapivě umírněnou
bilanci a~zároveň výrazný posun v~očekávané struktuře dotčených
profesí.\par

Weber v~makroekonomickém modelu
Q-INFORGE rozkládá dopad scénáře Industrie 4.0 v~horizontu do~roku~2030 na~německý
trh práce do~úbytku přibližně 490~tis.\ pracovních míst a~paralelního vzniku
zhruba 430~tis.\ nových pozic; čistý dopad ($\approx{-60}$~tis.) je~v~mezích
statistické chyby modelu. Klíčovým zjištěním není čistá bilance, ale~strukturální
přesun mezi odvětvími a~kvalifikačními skupinami: rostoucí přidaná hodnota
plyne k~vysoce kvalifikovaným profesím a~ke~komplementárním službám, zatímco
úbytek dopadá na~rutinní výrobní pozice. Weber explicitně poukazuje
na~nutnost právního a~smluvního rámce, který tento přesun zprostředkuje
prostřednictvím podnikové participace (\emph{Mitbestimmung}) a~kolektivně
vyjednaných rekvalifikačních fondů; bez nich hrozí, že~se~přínosy automatizace
zkoncentrují u~kapitálu a~náklady přizpůsobení dopadnou výhradně na~zaměstnance.~\cite{Weber_Industrie40_2015}\par

Zkušenosti posledních let posunuly původní očekávání ohledně toho, které
profese budou zasaženy nejdříve. Zatímco scénáře Industrie 4.0
předpokládaly rychlou erozi rutinní fyzické manipulace v~tovární výrobě,
empirické případy ukazují, že~mechanická automatizace v~nestandardizovaném
prostředí naráží na~tvrdé technické limity: projekt \emph{Fuselage Automated
Upright Build} společnosti Boeing (sestava trupů~777) byl po~čtyřech letech
provozu v~roce~2019 zrušen kvůli vyšší chybovosti než~u~ručního sestavení;
podobně zaváděná automatizace Tesly Model~3 byla v~roce~2018 částečně vrácena
k~lidským operátorům s~tím, že~„nadměrná automatizace byla
chyba”.\par

V~posledních třech letech je dopad zrychlujícího technologického pokroku patrný v~jiném směru ---~v~rutinní, nízkokvalifikované kognitivní práci. Analýza pro~Anthropic Economic Index kombinuje teoretickou expozici \acdat{LLM} se~skutečným využitím modelů
v~zákaznických datech a~odhadují průměrnou expozici
napříč profesemi v~USA na~zhruba \SI{30}{\percent}, s~výrazným rozptylem:
programátoři přibližně \SI{75}{\percent}, pracovníci zákaznického servisu
a~datového vstupu okolo \SI{67}{\percent}, řemeslné a~obslužné profese
prakticky nulová expozice. Přibližně \SI{30}{\percent} pracovníků zůstává
zcela mimo dosah současných \ac{LLM} systémů. Autoři rovněž upozorňují,
že~zhruba \SI{57}{\percent} pracovního využití generativní \acs{AI} připadá
na~\emph{augmentaci} (rozšíření schopností pracovníka), nikoli na~plnou
automatizaci úkolu ---~rozdíl, který má~přímé implikace pro~obsah \acgen{KV}:
augmentace nevede primárně k~náhradě pracovníka, ale k~úpravě jeho mzdové
struktury vůči zvýšené produktivitě. Studie zatím nenachází významný
dopad na~celkovou nezaměstnanost exponovaných skupin, ale identifikuje
mírný pokles náboru u~věkové skupiny 22--25~let v~exponovaných profesích
(řádově \SI{14}{\percent}), což naznačuje, že~první vlna dopadu může zasáhnout
generaci vstupující na~trh práce ---~právě tu, jejíž integrace tradičně
probíhala přes juniorní kancelářské pozice. To však představuje nejen krátkodobý
problém pro danou věkovou kohortu, ale v~dlouhodobém horizontu také výzvu pro
kontinuitu kvalifikované pracovní síly, pokud~se~neprovede včasná adaptace vzdělávacího systému a~sociálního dialogu.~\cite{Massenkoff_LaborMarketAI_2026}\par

Pro výklad těchto čísel je důležité nepovažovat současnou uživatelskou cenu generativní \acs{AI} za~stabilní měřítko dlouhodobé ceny práce. Nízká cena přístupu k~modelu může v~počáteční fázi fungovat jako penetrační, fakticky dumpingová cena digitální služby: zaměstnavatel neporovnává pouze cenu tokenu nebo licence, ale celý pracovní proces včetně integrace, kontroly kvality, odpovědnosti, datové bezpečnosti a~řízení chyb. Pokud se~\acs{AI} stane nezbytnou infrastrukturou pro kvalifikovanou kognitivní práci, její ekonomická cena se~může postupně přibližovat hodnotě práce, kterou nahrazuje nebo výrazně zesiluje. Pro~\ac{KV} z~toho plyne, že~předmětem vyjednávání nemá být jen prostý zákaz či povolení nástroje, ale transparentnost jeho nasazení, rozdělení produktivitních zisků a~ochrana vstupních pozic, na~nichž se~kvalifikace pracovníků teprve vytváří.~\cite{Massenkoff_LaborMarketAI_2026}~\cite{Petit_AIEmploymentImpact_EU_2026}\par

Evropská scénářová evidence tento závěr zpřesňuje pro~\acloc{ČR}. Model Evropské komise pro~dopad \acs{AI} a~navazujících digitálních technologií neukazuje plošný zánik zaměstnanosti: pro~\acacc{EU} odhaduje technologický příspěvek k~zaměstnanosti do~roku~2030 na~+\SI{3.7}{\percent} oproti kontrafaktuálnímu scénáři a~pro~\acacc{ČR} na~+\SI{6.2}{\percent}. Rozpad podle skupin je však podstatnější než agregát. U~mladých pracovníků ve~věku 15--24~let vychází pro~\acacc{ČR} do~roku~2030 technologický efekt přibližně -\SI{17.9}{\percent}, zatímco u~vysokokvalifikovaných pracovníků +\SI{18.9}{\percent}. Hlavním rizikem tedy není samotná automatizace, nýbrž nerovnoměrné rozdělení jejích přínosů a~nákladů: generace vstupující na~trh práce a~nízkokvalifikované profese potřebují přechodové mechanismy \aclgen{APZ} a~\aclgen{KV}, jinak se~pozitivní agregátní efekt může spojit s~hlubší mzdovou a~kariérní polarizací.~\cite{Petit_AIEmploymentImpact_EU_2026}\par

\ac{ČR} se~k~očekávané vlně strukturálních změn přibližuje s~jedním
z~nejnižších adaptačních rozpočtů v~\acs{OECD}: výdaje na~\acdat{APZ}
(viz~obr.~\ref{fig:eu_apz_vydaje_2004}) dlouhodobě nepřekračují
\SI{0,3}{\percent}~\acgen{HDP}, oproti dánským přibližně \SI{2}{\percent}
a~průměru \acs{geo-EU27} kolem \SI{0,5}{\percent}.
Slabinou není pouze absolutní velikost rozpočtu, ale zejména jeho alokační
profil: meta-analýzy \acs{OECD} dokládají nejvyšší dlouhodobou návratnost
u~cílených individualizovaných rekvalifikačních programů a~veřejného
zaměstnávání s~mentoringem, zatímco krátkodobá veřejně prospěšná opatření bez
vzdělávací složky vykazují nejnižší efekt.
Český systém aktuálně inklinuje spíše k~druhé kategorii opatření. Pokud má~\acgen{ČR}
absorbovat dopad automatizace bez prohloubení mzdové polarizace,
vyžaduje paralelní posun:
\begin{itemize}[noitemsep, topsep=2pt, leftmargin=*]
	\item navýšení podílu nákladů na \acgen{APZ} na~vzdělávání a~rekvalifikaci,
	\item vznik bipartitních sektorových fondů financovaných na základě dohody sociálních partnerů
	po~vzoru německého Bildungsgutscheinu nebo nizozemských \emph{O\&O}-fondů,
	\item zařazení doložek o~rekvalifikaci a~přerozdělení zisků ze~zvýšené produktivity do~\acs{KS}.
\end{itemize}
Bez koordinovaného sociálního dialogu, který by~tyto mechanismy ukotvil,
hrozí, že~strukturální fluktuace dopadne primárně na~jednotlivce
prostřednictvím nezaměstnanosti a~pomalého reálného mzdového růstu.~\cite{OECD_GoodJobsForAll_2018}~\cite{NERV_NavrhyRustu2025}~\cite{InovacniDoporuceni_ZahraniceKV}\par
